\section{Lista práctica de Recipe: cómo convertir los datos en capacidades}

Data Recipe no es un eslogan, sino un proceso experimental. Si un equipo afirma tener una ventaja en datos, al menos debe poder responder: de dónde vienen los datos, cómo se filtran, cómo se mezclan, cómo entran al entrenamiento, cómo se evalúan, cómo regresan los fallos y cómo se ajusta la siguiente ronda.

\subsection{Lista práctica de siete pasos}

\noindent\begin{longtable}{p{0.16\textwidth}p{0.36\textwidth}p{0.39\textwidth}}
\toprule
Paso & Acción & Punto de verificación \\
\midrule
1. Definir la capacidad & Precisar qué capacidad de tarea se quiere mejorar & ¿Es percepción, planificación, acción, recuperación o generalización? \\
2. Buscar fuentes de datos & Reales, simulados, sintéticos, de internet, retroalimentación humana & ¿Los datos cubren la capacidad objetivo? \\
3. Limpiar y filtrar & Eliminar duplicados, ruido, contenido irrelevante y contaminación & ¿Se eliminaron por error casos de cola larga y casos de fallo? \\
4. Fijar las proporciones & Determinar la proporción de cada fuente y de cada nivel de dificultad & ¿Provoca un desequilibrio entre capacidades? \\
5. Currículo de entrenamiento & Decidir qué se aprende primero y qué después & ¿Respeta la progresión de las capacidades? \\
6. Evaluación y retroalimentación & Verificar con benchmark, simulación o tareas reales & ¿Las métricas se relacionan con el despliegue real? \\
7. Ciclo cerrado de retorno & Devolver los fallos, la cola larga y las tareas nuevas al conjunto de datos & ¿Los datos de la siguiente ronda tienen mayor rendimiento marginal? \\
\bottomrule
\end{longtable}

\begin{knowledgebox}{Lo que se puede replicar de una Recipe y lo que no}
El proceso se puede replicar, pero la receta concreta es difícil de replicar, porque surge de la interacción entre el modelo, la tarea, el hardware, la evaluación y la experiencia del equipo; si cambia la capacidad objetivo, quizá haya que rehacer la receta.
\end{knowledgebox}

\subsection{Resumen de la sección}

Para convertir los datos en capacidades hay que partir de la capacidad objetivo, no de los datos que se tienen a la mano. Recipe es ingeniería de datos orientada a capacidades.
